\begin{definition}[Chronological product of the endpoint operations]%
    \label{def:peps_torus_swept_physical_route_composition}
    \lean{TNLean.PEPS.torusSweptPhysicalRouteProduct}
    \leanok
    \uses{def:peps_torus_swept_route_bond_assignments}
    For twelve physical matrices $M_0,\ldots,M_{11}$, let $P_0=I$ and
    $P_{n+1}=M_nP_n$ for $0\leq n<12$. Thus the latest operation acts
    on the left. This is the chronological product for the explicit
    endpoint route of \cite[\S6.6, local lines 2340--2423]{Schuch2010PEPS}.
\end{definition}

\begin{theorem}[Twelve fixed original-spin endpoint operations]%
    \label{thm:peps_torus_swept_physical_route_composition}
    \lean{TNLean.PEPS.IsGIsometric.exists_unitary_torusSweptPhysicalRoute}
    \leanok
    \uses{def:peps_torus_swept_physical_route_composition,
        thm:peps_torus_swept_route_vacancies,
        thm:peps_vacant_plaquette_physical_move}
    Assume torus periods at least eight and seven and regular
    G-isometry of every original site tensor. There are twelve global
    unitaries $M_0,\ldots,M_{11}$, chosen before both initial flux labels.
    Each is the identity extension of a unitary on its actual six-site
    movement patch, or the adjoint of such an extension.
    If $\Psi_n(g,h)$ is the actual closed tensor contraction with the
    recursively derived bond assignment after $n$ steps, then
    \begin{align}
        M_n\Psi_n(g,h) & =\Psi_{n+1}(g,h) &  & (0\leq n<12),    \\
        P_n\Psi_0(g,h) & =\Psi_n(g,h)     &  & (0\leq n\leq12).
    \end{align}
    Every $P_n$ is unitary. No flux label is used to choose these
    original-spin operations. This is the specified finite-torus
    endpoint route, not an unrestricted braid or an operation
    annihilating the final partners.
\end{theorem}

\begin{proof}\leanok
    At each source position choose the fixed elementary original-spin
    unitary in the required direction. The actual next plaquette has
    identity holonomy by the derived vacancy calculation, so its
    all-input action gives the literal next bond assignment. Choose all
    twelve operations before introducing the flux labels. Induction
    on the chronological product proves both its unitarity and the
    actual closed-state identity at every stage.
\end{proof}
